\section{The Batched Ballot: Example, Comparison, and Proof Sketches}\label{app:batched}
\textbf{Worked example:} Take the ballot of Table~\ref{tab:example-ballot}, $s = (0, 1, 1)$, so
$m = 3$ and $L = 2$ (values of $v$ below $4$). Cell $(A, B)$ has $\delta
= 1$, $b = 1$, $p = 1$, $v = -1 + 2 - 1 = 0$, bits $X = (0, 0)$. Cell
$(B, A)$ has $\delta = -1$, $b = 0$, $p = 0$, $v = 1$, $X = (1, 0)$. The
tie $(B, C)$ has $\delta = 0$, $b = 0$, $p = 0$, $v = 0$. In the two
rejected ballots of the table, claiming $b_{BA} = 1$ with $\delta_{BA} =
-1$ forces $p = -1$ and $v = 1 - 2 - 1 = -2$, and claiming $b_{BC} = 1$
with $\delta_{BC} = 0$ forces $v = -1$. Neither value is a sum of two
bits, so no $X$ satisfies (V3) and the ballot fails. The binary form of
$v$ replaces the $2m+1$ membership branches: a cheating voter is caught
by an equation with no solution instead of by a disjunction with no true
branch.

\begin{table}[h]
  \caption{The two ballots for $m$ candidates, $L = \lceil \log_2 (m+1)
  \rceil$. Measured values are in Table~\ref{tab:ballotcost}.}
  \label{tab:twoballots}
  \centering\scriptsize\setlength{\tabcolsep}{3pt}
  \begin{tabularx}{\columnwidth}{@{}l>{\raggedright\arraybackslash}X>{\raggedright\arraybackslash}X@{}}
    \toprule
     & first ballot (\S\ref{sec:ballot}) & batched ballot (\S\ref{sec:batched}) \\
    \midrule
    ciphertexts & $m + m^2$ & $m^2$ \\
    commitments & --- & $L + 3$, one key of $m^2$ generators \\
    proofs & $m$ membership, $m^2$ 0/1, $m^2$ membership & $1$ linear, $1$ zero argument of $L + 3$ rows \\
    size & $O(m^3)$ group elements & $2m^2 + O(\log m)$ group, $3m^2 + O(m)$ field elements \\
    prover, verifier & $O(m^3)$ exp. & $O(m^2 \log m)$, $O(m^2)$ exp. \\
    oracle queries & $1$ & $5$ \\
    soundness & exact, $2$ transcripts & statistical, $O(m^2 \log^2 m)$ transcripts \\
    group hyp. & $g$ of order $q$, char.\ $> 4m+2$ & char.\ $> 3 \cdot 2^L + 2$ \\
    key hyp. & none & binding of the commitment key \\
    zero-knowledge & perfect, no hypothesis & perfect, key in $\langle h_1 \rangle$ \\
    Rocq & \rocq{ballot_sound}, \rocq{ballot_zk} & \rocq{ballot3_sound}, \rocq{ballot3_zk} \\
    \bottomrule
  \end{tabularx}
\end{table}

\textbf{Soundness (\rocq{ballot3_sound}).} The proof runs the extractors in sequence. From two challenges of (a) at
the base $\beta$, \rocq{lin_open3} gives a witness $w$ with $\varphi(w) =
T$, hence $s$ with $C_\delta = \mathrm{com}(Ms; \rho_\delta)$ by
\rocq{com_gd}, an opening $b$ of $C_b$, the link $C_\beta =
\mathrm{Enc}(\langle \beta, b\rangle; R)$, and $C_0 = h_1^{\rho^*}$
(\rocq{phi3_target_open}). The vectors that vary one coordinate of
$\beta$ give further witnesses. By binding they open $C_b$ to the same
$b$, so the link holds for the same $b$ at every $\beta$. Dividing the
links at $\beta$ and at $\beta'$ differing at $k$ gives $c_k^{\beta_k -
\beta'_k} = \mathrm{Enc}(b_k (\beta_k - \beta'_k))$ and, taking the
$(\beta_k - \beta'_k)$-th root in the group, which is a vector space over
$\mathbb{Z}_q$, $\mathrm{Dec}(c_k) = g^{b_k}$ (\rocq{beta_bits},
\rocq{root_unique}), so the committed $b$ is the plaintext vector. The
ciphertext link is undone by roots and not by comparing exponents,
so the batched ballot needs no hypothesis on the order of $g$, while
the first ballot needs \rocq{exp_injective}. One accepted zero argument
gives, by the library's extractor, openings $p$ and $X^{(k)}$ of $C_p$
and the $C_{X^{(k)}}$ (\rocq{zero_open_ref}), and every other accepted
$(y, z)$ opens the same vectors unless binding is broken
(\rocq{zero_open_at}). The weighted sum is then $z$ times a polynomial
of degree $L + 1$ in $z$ whose coefficients are the row relations, so $L
+ 2$ values of $z$ make each row vanish (\rocq{coeffs_vanish}), and
$m^2$ values of $y$ make each row vanish coordinate-wise by the
library's Vandermonde argument (\rocq{prod_of_star},
\rocq{bits_of_star}). Hence $p = b \circ \delta$, and $b$ and the
$X^{(k)}$ are bit vectors. By binding on $C_0$, the committed vector
$\delta + b - 2p + \sum_k 2^k X^{(k)}$ is zero entry by entry
(\rocq{C0_open}). Finally \rocq{decode_positions} reads off the
positions. Each $v_k = \sum_t 2^t X^{(t)}_k$ is a natural number below
$2^L$. If $b_k = 0$ then $\delta_k = -v_k \le 0$, and if $b_k = 1$ then
$\delta_k = v_k + 1 \ge 1$, so $b_k = [\delta_k \ge 1]$. Since
$\delta_{ij} = s_j - s_i$, setting $\mathit{pos}_i = 2^L + (s_i - s_0)$,
a natural number in $[0, 2^{L+1}]$ by the same bounds, gives $b_{ij} =
[\mathit{pos}_i < \mathit{pos}_j]$. The characteristic bound lets these
integer inequalities be read in the field.

\textbf{Zero-knowledge (\rocq{ballot3_zk}).} The proof composes three bijections on the randomness: the library's
\rocq{lin_shvzk} and \rocq{zero_shvzk} for the proofs, given the
commitments, and $\rho \mapsto \rho + \langle \log_{h_1} \mathbf{g},
v\rangle$ for the commitments, which makes a commitment to $v$ uniform
in $\langle h_1 \rangle$. Reading $C_\delta$ as a commitment to $\delta$
under $\mathbf{g}$ means the derived key needs no logarithms of its own.
Like the library's Bayer--Groth proofs, the last step writes each
commitment generator as a power of $h_1$, $g_j = h_1^{e_j}$. Such
exponents exist for every key in the subgroup; the proof uses them only
to exhibit the bijection, and the simulator never needs them.
